For an explicit definition of every coefficient in the ensuing
linear system, retain the complete original recurrence
\cite[equations (4.1)--(4.2), (5.1)--(5.6)]{cqNS2026}.
In the present recurrence $A,D$ are the source's scalar
exponents, and $D_X$ is its logarithmic radial derivative:
\begin{equation}
\begin{gathered}
 A=\tfrac12+h,\qquad D=\tfrac12-h,\qquad
 \lambda_n=2nh,\qquad b=-A-\tfrac12,\qquad c=-A,\\
 d=1-\eta^2,\qquad L=1-2h\eta^2,\qquad D_X=X\partial_X,\\
 \tau=1-t=q(1-\eta^2),\qquad z=q^D\eta,\qquad
 \sigma=\frac{r^2}{2},\qquad X=\frac{\sigma}{q}.
\end{gathered}
 \label{eq:cq-source-recurrence-parameters}
\end{equation}
The source denotes the radial $\sigma$ in this line by $s$;
the equation explicitly identifies that source variable with
our retained radial coordinate, while leaving the Connes
spectral $s$ distinct. Every other source coefficient in the
recurrence is unchanged. For a real exponent $a$ define
\begin{align}
 T_a f&=L^{-1}(-af+D\eta\,\partial_\eta f+D_Xf),\notag\\
 Z_a f&=L^{-1}(2a\eta f+d\,\partial_\eta f-2\eta D_Xf),\notag\\
 T_{a,n}&=T_{a+\lambda_n},\qquad
 Z_{a,n}=Z_{a+\lambda_n},\qquad
 Z^{[2]}_{a,n}=Z_{a+\lambda_n-D}Z_{a+\lambda_n}.
 \label{eq:cq-source-profile-time-axial-operators}
\end{align}
The rightmost operator acts first in the last expression.
These are the exact physical derivative identities
\[
 \partial_t(q^af)=q^{a-1}T_af,\qquad
 \partial_z(q^af)=q^{a-D}Z_af.
\]
Indeed implicit differentiation of the unchanged similarity
coordinates gives
\[
 q_t=-L^{-1},\quad \eta_t=\frac{D\eta}{qL},\quad
 X_t=\frac{X}{qL},\qquad
 q_z=\frac{2\eta q^{1-D}}L,\quad
 \eta_z=\frac{d}{q^DL},\quad
 X_z=-\frac{2\eta X}{q^DL};
\]
substitution in the ordinary chain rule proves both formulas.
A second axial derivative changes the exponent by $-D$,
which proves the stated composition order for $Z^{[2]}$.

At each nonnegative order the physical profiles are
\begin{equation}
\begin{aligned}
 u_{\theta,n}
  &=q^{-A+\lambda_n}E_n
    =r q^{-A-1/2+\lambda_n}\phi_n/C,
 &E_n&=\sqrt{2X}\phi_n/C,\\
 u_{z,n}&=q^{-A+\lambda_n}U_n,
 &r u_{r,n}&=q^{\lambda_n}V_n,\\
 p_n&=q^{-2A+\lambda_n}\Pi_n,
 &A_Xf&=\int_0^1 f(b_0X,\eta)\,db_0 .
\end{aligned}
 \label{eq:cq-source-physical-coefficients}
\end{equation}
The letter $b_0$ in the average is an integration variable,
distinct from the exponent $b=-A-1/2$.
The leading profiles have index zero; every negative-index
field is zero. Primes in the following recurrences mean
$\partial_X$ at fixed $\eta$. Incompressibility gives
\begin{equation}
 \partial_XV_n=-Z_{-A,n}U_n,\qquad
 \frac{V_n}{X}=
 \frac{2\eta U_n-2\eta(D+\lambda_n)A_X(U_n)
                      -d\,\partial_\eta A_X(U_n)}L .
 \label{eq:cq-source-full-incompressibility-recurrence}
\end{equation}
The original positive-order equations are, for $n\geq1$,
\begin{align}
 2(X\phi_n''+2\phi_n')
 &=T_{b,n}\phi_n+
   \sum_{\substack{i,j\geq0\\i+j=n}}
    \left\{V_i(\phi_j'+\phi_j/X)+U_iZ_{b,j}\phi_j\right\}
                         -Z^{[2]}_{b,n-1}\phi_{n-1},
 \label{eq:cq-source-full-angular-recurrence}\\
 2(XU_n''+U_n')
 &=T_{c,n}U_n+
   \sum_{\substack{i,j\geq0\\i+j=n}}
           \left\{V_iU_j'+U_iZ_{c,j}U_j\right\}
   +Z_{-2A,n}\Pi_n-Z^{[2]}_{c,n-1}U_{n-1},
 \label{eq:cq-source-full-axial-recurrence}\\
 \Pi_n'
 &=C^{-2}\sum_{\substack{i,j\geq0\\i+j=n}}\phi_i\phi_j
                          -\frac{\Omega_{n-1}}{2X},
 \label{eq:cq-source-full-pressure-recurrence}\\
 \Omega_k
 &=T_{0,k}V_k+
   \sum_{\substack{i,j\geq0\\i+j=k}}
       \left\{V_i(V_j'-V_j/(2X))+U_iZ_{0,j}V_j\right\}
             -2XV_k''-Z^{[2]}_{0,k-1}V_{k-1}.
 \label{eq:cq-source-full-radial-recurrence}
\end{align}
In particular the pressure term is
$-\Omega_{n-1}/(2X)$, with the factor $2$ in its denominator.
These formulas define each lower-order source without leaving
a pressure, radial-viscous or axial-viscous contribution
implicit. All ordered pairs in the sums are retained.

The explicit six-row reduction can also be stated without an
unnamed source matrix. At fixed $n\geq1$ define the three known
lower-order functions
\begin{align*}
 f_n^\phi
 &=\sum_{i=1}^{n-1}
     \{V_i(\phi_{n-i}'+\phi_{n-i}/X)
                         +U_iZ_{b,n-i}\phi_{n-i}\}
                   -Z^{[2]}_{b,n-1}\phi_{n-1},\\
 f_n^U
 &=\sum_{i=1}^{n-1}
     \{V_iU_{n-i}'+U_iZ_{c,n-i}U_{n-i}\}
                   -Z^{[2]}_{c,n-1}U_{n-1},\\
 f_n^\Pi
 &=C^{-2}\sum_{i=1}^{n-1}\phi_i\phi_{n-i}
                                     -\Omega_{n-1}/(2X).
\end{align*}
Empty sums are zero. Put $\xi=\sqrt X$, $v_0=V_0/X$, and
write the six current unknowns as
\[
 W_n=(\phi,U,K,\Pi,p_\xi,q_\xi)^{\mathsf T},\qquad
 p_\xi=\partial_\xi\phi,\quad q_\xi=\partial_\xi U,\quad
 K=A_X(U)-U .
\]
Here $p_\xi,q_\xi$ are labels for the last two components;
they are not the physical pressure or the similarity
coordinate $q$. The current radial-velocity coefficient is
the following explicit linear expression:
\[
 \mathcal V_n[U,K]=
 \frac{2\eta U-2\eta(D+\lambda_n)(U+K)
                       -d\,\partial_\eta(U+K)}L .
\]
Thus $V_n=X\mathcal V_n[U,K]$. With every known coefficient
evaluated at $(X,\eta)=(\xi^2,\eta)$, the first four rows are
\begin{equation}
\begin{aligned}
 \partial_\xi\phi&=p_\xi,&
 \partial_\xi U&=q_\xi,\\
 \partial_\xi K+\frac2\xi K&=-q_\xi,&
 \partial_\xi\Pi
      &=4\xi C^{-2}\phi_0\phi+2\xi f_n^\Pi .
\end{aligned}
 \label{eq:cq-source-explicit-first-four-rows}
\end{equation}
The complete last two rows are
\begin{equation}
\begin{split}
 \partial_\xi p_\xi+\frac3\xi p_\xi
 =2\biggl\{&
 \frac{-(b+\lambda_n)\phi+D\eta\,\partial_\eta\phi
                                     +\xi p_\xi/2}{L}\\
 &+v_0(\xi p_\xi/2+\phi)
 +U_0\frac{2(b+\lambda_n)\eta\phi+d\,\partial_\eta\phi
                                           -\eta\xi p_\xi}{L}\\
 &+\mathcal V_n[U,K](X\phi_0'+\phi_0)
                    +U Z_{b,0}\phi_0+f_n^\phi\biggr\}.
\end{split}
 \label{eq:cq-source-explicit-angular-row}
\end{equation}
\begin{equation}
\begin{split}
 \partial_\xi q_\xi+\frac1\xi q_\xi
 =2\biggl\{&
 \frac{-(c+\lambda_n)U+D\eta\,\partial_\eta U
                                     +\xi q_\xi/2}{L}
                       +v_0\xi q_\xi/2\\
 &+U_0\frac{2(c+\lambda_n)\eta U+d\,\partial_\eta U
                                           -\eta\xi q_\xi}{L}
 +X\mathcal V_n[U,K]U_0'+U Z_{c,0}U_0\\
 &+\frac{2(-2A+\lambda_n)\eta\Pi+d\,\partial_\eta\Pi
              -2\eta X(2C^{-2}\phi_0\phi+f_n^\Pi)}L
                                      +f_n^U\biggr\}.
\end{split}
 \label{eq:cq-source-explicit-axial-row}
\end{equation}
These rows follow directly from
\eqref{eq:cq-source-full-incompressibility-recurrence}--%
\eqref{eq:cq-source-full-radial-recurrence}.
For the radial conversion,
$\partial_X\phi=p_\xi/(2\xi)$ and
$2(X\phi''+2\phi')=(\partial_\xi p_\xi+3p_\xi/\xi)/2$;
the analogous axial identity has $q_\xi/\xi$ instead of
$3p_\xi/\xi$. The apparent divisions in the transport rows
cancel exactly with $V_0=Xv_0$ and
$V_n=X\mathcal V_n[U,K]$, producing the displayed regular
coefficients. In the pressure contribution to the last row,
the substitution is the full equation
$\Pi'=2C^{-2}\phi_0\phi+f_n^\Pi$; hence its lower-order
$f_n^\Pi$ contribution is included.

Equations \eqref{eq:cq-source-explicit-first-four-rows}--%
\eqref{eq:cq-source-explicit-axial-row} define $A_0W_n+
A_1\partial_\eta W_n+f_n$ uniquely by collecting, in this
specified order of the six components, the coefficients of
$W_n$, $\partial_\eta W_n$, and the remaining known terms.
Thus the Picard argument's coefficient functions are
completely specified. Extracting the $\partial_\eta$ terms
from these rows gives exactly the sparse $A_1$ displayed
in that argument: the coefficient of $\partial_\eta\phi$
is $2(D\eta+dU_0)/L$, the $\partial_\eta U$ and
$\partial_\eta K$ coefficients in the angular row are both
$-2d(\phi_0+X\phi_0')/L$, and the axial row gives
$2(D\eta+dU_0-dXU_0')/L$, $-2dXU_0'/L$, and $2d/L$.
This supplies a direct check of every factor in the sparse
matrix and preserves every lower-order forcing term.
